\chapter{表达式、定义、函数与计算}

\section{每个表达式都有类型}

Lean 的核心判断写成
\[
\Gamma \vdash t : T,
\]
表示在局部环境 $\Gamma$ 中，项 $t$ 具有类型 $T$。例如自然数 \code{3} 的类型是 \code{Nat}，函数 \code{Nat.succ} 的类型是 \code{Nat -> Nat}，而一个等式证明也是某个命题类型的项。

\begin{lstlisting}
#check 3
#check Nat.succ
#check fun n : Nat => n + 1
#check Type
#check Prop
\end{lstlisting}

类型本身也需要类型。Lean 使用 universe 层级避免“所有类型组成的类型包含自身”带来的悖论。入门时可以先记住：\code{Nat : Type}，而 \code{Type : Type 1}；多态定义经常让 universe 参数由 elaborator 自动补全。

\section{函数类型与函数定义}

若输入类型为 $A$、输出类型为 $B$，函数类型写作 $A\to B$。匿名函数和命名定义分别写作：
\begin{lstlisting}
#check fun n : Nat => n * n

def square (n : Nat) : Nat := n * n

def compose (f : B -> C) (g : A -> B) : A -> C :=
  fun x => f (g x)
\end{lstlisting}
最后一个定义中的 \code{A B C} 需要预先作为类型变量声明：
\begin{lstlisting}
universe u v w
variable {A : Type u} {B : Type v} {C : Type w}
\end{lstlisting}

函数默认柯里化。\code{Nat -> Nat -> Nat} 解析为 \code{Nat -> (Nat -> Nat)}：先接受一个自然数，再返回一个接受自然数的函数。因此多参数应用写作 \code{f x y}，而不是必须构造元组。

\section{显式参数、隐式参数与命名参数}

圆括号参数必须显式提供，花括号参数通常由上下文推断：
\begin{lstlisting}
def identity {A : Type} (x : A) : A := x

#check identity 5
#check identity "lean"
#check @identity
\end{lstlisting}
\code{@identity} 要求显示包括隐式参数在内的完整函数。诊断类型推断时，还可以用命名参数指定：\code{identity (A := Nat) 5}。

\begin{warningbox}
“参数没有写出来”不等于“参数不存在”。mathlib 定理常含 universe、类型、实例和假设等多层隐式参数。应用失败时先 \code{\#check} 定理，再必要时查看 \code{@theoremName} 或打开更详细的 pretty printer。
\end{warningbox}

\section{局部定义、命名空间与模块}

\code{let} 创建局部名称；\code{namespace} 避免全局命名冲突：
\begin{lstlisting}
namespace SelfStudy

def twice (f : Nat -> Nat) (x : Nat) : Nat :=
  let y := f x
  f y

#eval twice (fun n => n + 1) 10

end SelfStudy
\end{lstlisting}
跨文件使用定义时，需要让文件路径、模块名与 \code{import} 对应。\code{open} 只改变名称访问方式，并不会导入尚未加载的模块。

\section{模式匹配与递归}

代数数据类型通过构造子产生值，模式匹配负责拆解值：
\begin{lstlisting}
def isZero : Nat -> Bool
  | 0 => true
  | _ + 1 => false

def sum : List Nat -> Nat
  | [] => 0
  | x :: xs => x + sum xs

#eval sum [2, 3, 5]
\end{lstlisting}
\code{sum} 的递归调用作用于严格更小的尾列表，所以 Lean 的终止检查器可以接受。逻辑一致性要求一般递归不能无条件地产生任意类型的值；若递归不是结构递减，需要提供终止度量与证明。

\begin{intuitionbox}
函数式定义同时承担三种角色：它是可执行程序，是数学对象，也是后续证明展开的重写规则。定义接口若写得清楚，计算和证明会共享同一份结构；定义若混入过多实现细节，后续 theorem 也会被这些细节绑住。
\end{intuitionbox}

\section{归约与定义性相等}

Lean 中两类“相等”需要分开：
\begin{description}
  \item[定义性相等] 通过展开定义、函数应用归约或模式匹配计算后相同，kernel 可直接识别；
  \item[命题相等] 类型为 \code{Eq a b} 的命题，需要一个证明项。
\end{description}

例如：
\begin{lstlisting}
def addOne (n : Nat) := n + 1

example : addOne 2 = 3 := by
  rfl

example (n : Nat) : n + 0 = n := by
  simpa using Nat.add_zero n
\end{lstlisting}
第一条可通过计算完成；第二条对变量 $n$ 并不是简单把两边算成数字，需要自然数加法的定理（该定理本身由归纳证明）。

\section{可执行不等于已证明}

\code{\#eval} 能展示某个输入的计算结果，但它没有证明对所有输入成立。例如运行排序函数得到有序列表，只验证了一个样例。一般性正确性需要分别说明：输出有序、输出是输入的排列，以及函数终止。

\begin{aibox}
AI4TCS 中模型可以生成候选程序，\code{\#eval}、单元测试和随机测试用于快速筛选；Lean theorem 用于证明明确写入类型的性质。测试与证明共享定义，却提供不同等级的证据，不能互相冒充。
\end{aibox}

\section{定义设计检查表}

在开始证明前，先检查：
\begin{enumerate}
  \item 输入、输出和错误情况是否在类型中表达清楚；
  \item 使用 \code{Nat}、\code{Int}、\code{Fin n} 还是一般类型是否符合问题；
  \item 无效状态能否被构造，若能，后续需要什么不变量；
  \item 定义是面向规范还是面向高性能实现，两者是否需要分层；
  \item 递归为何终止，模式是否覆盖所有构造子。
\end{enumerate}

\begin{checkpointbox}
\begin{enumerate}
  \item 用 \code{\#check} 解释五个表达式的完整类型；
  \item 写出一个带隐式类型参数的多态函数，并用 \code{@} 查看完整参数；
  \item 定义列表乘积或最大值的安全版本，说明空列表怎样处理；
  \item 区分定义性相等与需要证明的命题相等；
  \item 解释为什么 \code{\#eval} 的十个样例不能代替全称定理。
\end{enumerate}
\end{checkpointbox}
